\documentclass[11pt,compress,t,notes=noshow, aspectratio=169, xcolor=table]{beamer}

\usepackage{../../style/lmu-lecture}
% Defines macros and environments
\input{../../style/common.tex}

\title{Applied Machine Learning}
% \author{LMU}
%\institute{\href{https://compstat-lmu.github.io/lecture_iml/}{compstat-lmu.github.io/lecture\_iml}}
\date{}

\begin{document}

% https://sebastianraschka.com/pdf/lecture-notes/stat479fs18/11_eval-algo_slides.pdf

% Raschka paper
% https://arxiv.org/pdf/1811.12808.pdf
\begin{frame}{Hypothesis Testing in Benchmarking}

\begin{itemize}
    \setlength\itemsep{2em}
    \item We want to ensure that a model consistently outperforms another model, not just on certain test sets.
    \item Hypothesis tests quantify such evidence, meaning that our benchmark describes the general behavior of the algorithms and not a coincidental finding specific to a certain test set.
    \item We can think of 3 different test scenarios:
        \begin{itemize}
            \item 2 models / algorithms on a single domain (i.e., on a single data set)
            \item 2 algorithms across different domains (i.e., on multiple data sets)
            \item multiple algorithms across different domains / data sets.
        \end{itemize}
    \item In this case, a data set consists of multiple performance measurements of each model.
    \end{itemize}

    
\end{frame}

% \begin{frame}{Hypothesis Testing - General Procedure}

% \begin{itemize}
%     \setlength\itemsep{2em}
%     \item Formulate hypothesis, e.g., for a two-sided test.
    
%     $H_0$: The model performance of two models is identical.
    
%     $H_1$: The model performance of two models is different.
    
%     \item Given distribution of test statistic, compute threshold value for given significance level $\alpha$.
%     \item Compute test statistic and p-value given validation data set.
%     \item If p-value < $\alpha$: Reject $H_0$ (the model performance is different).
% \end{itemize}

    
% \end{frame}

% \begin{frame}{Hypothesis Tests (not a slide)}

% \begin{itemize}
%     \setlength\itemsep{2em}
%     \item McNemar's test to evaluate classifiers.
%     %https://rasbt.github.io/mlxtend/user_guide/evaluate/mcnemar/
%     \item Difference in proportions, by Snedecor and Cochran
%     \item Resampled paired t-test
%     \item k-fold cross-validated t-test
%     \item 5x2cv paired t-test
%     \item F-test for classifiers
% \end{itemize}
    
% \end{frame}

% \begin{frame}{McNemar's Test for Classifiers}

% \begin{itemize}
%     % \setlength\itemsep{1em}
%     \item McNemar's test compares the errors of two models in a 2x2 contingency matrix:
%     \begin{figure}
%         \centering
%         \includegraphics[width = 0.6\textwidth]{slides/04-perf-eval/mcnemar_contingency_table_ex1.png}
%         \label{fig:enter-label}
%     \end{figure}
%     \item Both models seem to perform equally well (model 1 accuracy: 9,960 / 10,000 = 99.6\%; model 2 accuracy: 9,970 / 10,000 = 99.7\%).
%     \item In subfigure A, model 2 got 11 predictions right that model 1 got wrong; model 2 got 1 prediction right that model 2 got wrong (ratio 11:1).
%     $\Rightarrow$ Model 2 seems to perform better than model 1.
%     \item In subfigure B, the ratio is 25:15, which makes it less clear to choose a better-performing model.
% \end{itemize}
    
% \end{frame}


% \begin{frame}{McNemar's Test for Classifiers}

% \begin{itemize}
%     % \setlength\itemsep{1em}
%     \item This relative error rate of two models can be used for a statistical test.
% \end{itemize}
    
% \end{frame}

% \begin{frame}{Benchmark Experiments}

% \begin{itemize}
% \item In benchmark experiments, different learning algorithms are applied to one or more datasets with the aim of comparing and ranking their performances.
% \item To ensure comparability, synchronized train and test sets, i.e., the same resampling method with the same train-test splits, should be used to calculate and compare the performances.
% \item Results of benchmark experiments produce a dataset which can be further analyzed and visualized. 
% \newline
% \textbf{Example}: Benchmark results (per CV-fold) of CART and random forest using 2-fold CV with MSE as performance measure: [...]
% \end{itemize}

% \end{frame}


% \begin{frame}{Hypothesis Testing in Benchmarking}
% We want to know if the difference in performance between models (or algorithms) is significant or if it only occurred by chance.

% \lz
% \textbf{Null Hypothesis Statistical Testing (NHST)} in benchmarking:

% \begin{itemize}
% \item Formulate a null hypothesis $H_0$ (e.g., the expected generalization error of two algorithms is equivalent).
% \item Use a hypothesis test to reject $H_0$.
% \item Rejecting $H_0$ gives us some confidence in the belief that the observed results may be not merely random.
% \end{itemize}

% % \begin{multicols}{2}
% % p-value: the probability that, when the null hypothesis is true, the statistical summary would be greater than or equal to the actual observed results
% %
% % Reject $H_0$ if $p<\alpha$
% %
% % \begin{flushright}\includegraphics[width=0.5\textwidth]{figure_man/table-of-errors.png} \end{flushright}
% % \end{multicols}
% % Typical example in medicine:
% %
% % \begin{itemize}
% % \item $H_0$: on average, the new drug is not better than the current drug
% % \item $H_a$: on average, the new drug is better
% % \end{itemize}

% \lz Typical example in machine learning:

% \begin{itemize}
% \item $H_0$: on average, model 1 does not perform better than model 2.
% \item $H_1$: on average, model 1 outperforms model 2.
% \item Aim: Reject $H_0$ with confidence level of $1-\alpha$.
% \end{itemize}
% \end{frame}

% \begin{frame}{Hypothesis Testing in Benchmarking}


% %It is important to clearly define what should be tested with $H_0$ and $H_1$, i.e., in the previous example, $H_0$ does not take into account \textit{how much} learner A outperforms learner B. With enough data even very small differences can be significant.
% There are 3 different scenarios:
% \begin{itemize}
% \item 2 models / algorithms on a single domain (i.e., on a single dataset),
% \item 2 algorithms across different domains (i.e., on multiple datasets),
% \item multiple algorithms across different domains / datasets.
% \end{itemize}

% Overview of different hypothesis tests for matched samples:
% \begin{center}\includegraphics[width=\textwidth]{slides/04-perf-eval/figure/tests_overview.png} \end{center}

% \footnotesize
% {\tiny{Source: \code{\url{https://marthawhite.github.io/mlcourse/lectures/2017/Lec20-MeasuringPerformance.pdf}}, p. 142.}\par}

    
% \end{frame}

% xx
% xx

% \begin{frame}{Two-matched-samples t-test}

% \begin{itemize}
%     \setlength\itemsep{1em}
%     \item The two-matched-samples t-test (i.e., a paired t-test) is the simplest hypothesis test if the aim is to compare two \textbf{models} on a single test set based on arbitrary performance measures.
%     \item Due to being a parametric test, several distributional assumptions must be made:
    
%     \begin{itemize}
%         \setlength\itemsep{1em}
%         \item \textbf{(Pseudo-)normality}, usually met when sample size > 30.
%         \item \textbf{i.i.d. samples}, usually met if the loss of individual observations from a single test set is considered. % (this assumption is violated in case of resampling as ). %difficult as data is often limited (assumption is violated in the case of resampling).
%         \item \textbf{Equal variances of populations}, can be investigated by plots.
%     \end{itemize}
% \end{itemize}



% \end{frame}

\begin{frame}{Two-matched-samples t-test}

    
    % A paired t-test to compare two different models $\fh_1$ and $\fh_2$ w.r.t. a performance measure calculated on a test set of size $n_{\text{test}}$:
    
    \begin{itemize}
        \setlength\itemsep{1em}
        \item $H_0$: $\widehat{\text{GE}}(\mathcal{I}_1, \mathcal{J}, \lamv_1, \rho) = \widehat{\text{GE}}(\mathcal{I}_2, \mathcal{J}, \lamv_2, \rho)$ \\ 
        $H_1$: $\widehat{\text{GE}}(\mathcal{I}_1, \mathcal{J}, \lamv_1, \rho) \neq \widehat{\text{GE}}(\mathcal{I}_2, \mathcal{J}, \lamv_2, \rho)$
        \item Test statistic $T = \sqrt{n} \frac{\bar{d}}{\sigma_{d}}$:
        %Teststatistic $t=\frac{d\left(\fh_A,\fh_B\right)}{\sqrt{\sigma_{A}^2 + \sigma_{B}^2}}$, with
        $$\sigma_{d} = \sqrt{\frac{1}{n - 1}\sum_{i=1}^{n} \left(d_i - \bar{d} \right)^2}
        $$
        where $d_i = L(\yi, \fh_1 (\xi)) - L(\yi, \fh_2 (\xi))$ and $\bar{d} = \sum\limits_{i=1}^{n} d_i$
        \item The two-matched samples t-test (or paired t-test) is one of the most popular tests for benchmarking but has multiple shortcomings:
        \begin{itemize}
            \item \textbf{Randomness:} Performance samples need to be random samples from the underlying population.
            \item \textbf{(Pseudo-) Normality:} The t-test requires a minimum of $\approx$ 30 samples or that samples come from a normal distribution.
            \item \textbf{Equal variances:} The paired t-test assumes equal variances of both populations.
            \end{itemize}
        \end{itemize}
\end{frame}
% xx

% \begin{frame}{Two-matched-samples t-test}
%     % \textbf{Note}: Here, $d_i$ is the difference of the outer loss of individual observations from the test set between the two models to be compared.
%     %\textbf{Issue}: $\GEh{n} \left(\fh_A\right)$ and $ \GEh{n} \left(\fh_B\right)$ are not independent as they are calculated on the same test set.
    
%     \begin{itemize}
%         \setlength\itemsep{1em}
%     \item We could also use a \textbf{$k$-fold CV paired t-test} to compare two \textbf{training algorithms} (instead of two models) on a single dataset.
%     \item Instead of comparing the outer loss of individual observations we would then compare the individual generalization errors per CV fold (i.e., the generalization error of the $k$ prediction models induced by the learning algorithm in each CV fold).
%     %\item Instad of using the difference of the outer loss of individual observations
%     \item Although the test sets do not overlap, the performance differences are not independent across CV folds due to overlapping training sets (which violates the assumption of i.i.d. samples).
%     %However, the paired t-test is a parametric test and is not recommended in practice as the assumptions of the student's t-test are often violated.
%     %\item A probably better alternative is, therefore, the Friedman test.
%     \item To partly overcome the issue of overlapping training sets across folds, Dietterich (1998) \footnote{Dietterich (1998). Approximate statistical tests for comparing supervised classification learning algorithms.}suggests using 5 times 2-fold CV so that at least within each repetition neither the training nor the test sets overlap.
%     % footnotes do not work
%     \end{itemize}
% \vspace{2cm} {\scriptsize Dietterich (1998). Approximate statistical tests for comparing supervised classification learning algorithms.}
% \end{frame}





% \begin{frame}{Weakness of Paired T-Test}

% \begin{itemize}
%     \setlength\itemsep{2em}
%     \item \textbf{Commensurability:} One has to be able to compare the differences over the data
%     sets, thus if the used performance measures have different properties for different
%     tasks, this assumption might be violated.
%     \item \textbf{(Pseudo-) Normality:} The t-test requires that (performance) samples come from a normal distribution or a minimum of 30 samples.
%     \item \textbf{Outliers:} Extreme values might skew the statistics and thus decrease a tests
%     power.
% \end{itemize}

% \end{frame}


\begin{frame}{McNemar Test for Classification}
% https://slideplayer.com/slide/9574350/
\begin{itemize}
\item The McNemar test compares the relative errors of two models.
\item As opposed to the paired t-test, it is non-parametric and does not rely on distributional assumptions.
\item Test statistic:
$$\chi^2_{Mc} =  \tfrac{(|B-C| - 1)^2}{B + C} \sim \chi^2_{1}
$$
% when the considered performance measure is based on an outer loss with a nominal or binary output, e.g., accuracy is based on a binary outer loss.
% \item Both models are trained on a training set and evaluated on a test set. Based on the test set, a \textbf{contingency table} that compares the two models (model 1 and model 2) is calculated:
\end{itemize}

\begin{minipage}{0.45\textwidth}
\includegraphics[width=\textwidth]{slides/04-perf-eval/figure/mcnemar_1.png} \end{minipage}
\begin{minipage}{0.54\textwidth}
\begin{itemize}
\item A: $\#$ obs. correctly classified by both.
\item B: $\#$ obs. misclassified by model 2 but not by model 1.
\item C: $\#$ obs. misclassified by model 1 but not by model 2.
\item D: $\#$ obs. misclassified by both.
\end{itemize}
\end{minipage}
\end{frame}

% \begin{frame}{McNemar Test}

% \begin{minipage}[c]{0.4\linewidth}
%   \begin{center}
%   \renewcommand{\arraystretch}{1.5}
%   \begin{tabular}{cc|cc}
%       & & \multicolumn{2}{c}{$f_2$} \\
%       & & $0$ & $1$ \\
%       \hline
%       \multirow{2}{*}{ $f_1$} & $0$ & $c_{00}^{Mc}$ & $c_{01}^{Mc}$ \\
%       & $1$ & $c_{10}^{Mc}$ & $c_{11}^{Mc}$ \\
%   \end{tabular}
%   \end{center}
% \end{minipage} % no space if you would like to put them side by side
% \begin{minipage}[c]{0.59\linewidth}
% \vspace{0.15cm}
%   $c_{00}^{Mc} = \sum_{i=1}^{|S_{test}|} \left[ I(f_1(\mathbf{x_i}) \neq y_i) \wedge I(f_2(\mathbf{x_i}) \neq y_i) \right]$ \vspace{0.05cm}

%   $c_{01}^{Mc} = \sum_{i=1}^{|S_{test}|} \left[ I(f_1(\mathbf{x_i}) \neq y_i) \wedge I(f_2(\mathbf{x_i}) = y_i)\right]$ \vspace{0.05cm}

%   $c_{10}^{Mc} = \sum_{i=1}^{|S_{test}|} \left[ I(f_1(\mathbf{x_i}) = y_i) \wedge I(f_2(\mathbf{x_i}) \neq y_i)\right]$ \vspace{0.05cm}

%   $c_{11}^{Mc} = \sum_{i=1}^{|S_{test}|} \left[ I(f_1(\mathbf{x_i}) = y_i) \wedge I(f_2(\mathbf{x_i}) = y_i)\right]$

% \end{minipage}
% \vspace{0.5cm}

% Testing: $H_0$: $c_{01}^{Mc} = c_{01}^{Mc} = c_{null}^{Mc}$

% $$\chi^2_{Mc} =  \frac{(|c_{01}^{Mc} - c_{01}^{Mc}| - 1)^2}{c_{01}^{Mc} + c_{10}^{Mc}} \overset{approx}{\sim} \chi^2_{1,1-\alpha}$$

% If $c_{01}^{Mc} + c_{01}^{Mc} < 20$, approximation doesn't work, use binomial instead
    
% \end{frame}

% \begin{frame}{McNemar Test}


% \begin{minipage}[c]{0.625\linewidth}
% Given such a contingency table, the accuracy of each model can be computed as follows:
% \begin{itemize}
%   \item Model 1: (A+B)/(A+B+C+D)
%   \item Model 2: (A+C)/(A+B+C+D)
% \end{itemize}

% Even if the models have \textbf{equal} accuracy (indicating equal performance), cells B and C may differ because the models may misclassify different instances.
% \end{minipage}
% % \begin{minipage}[c]{0.365\linewidth}
% %   \includegraphics{figure_man/mcnemar_1.png}
% % \end{minipage}

% \lz McNemar tests the following hypotheses:
% \begin{itemize}
% \item $H_0:$ Both models have the same performance (we expect B = C).
% \item $H_1:$ Performances of the two models are not equal.
% \end{itemize}
% %\lz
% The test statistic is computed by
% $$\chi^2_{Mc} =  \tfrac{(|B-C| - 1)^2}{B + C} \sim \chi^2_{1}.$$

% \textbf{Note}: The McNemar test should only be used if $B + C > 20$.
% \end{frame}

\begin{frame}{McNemar Test}

% The test statistic is computed by
% $$\chi^2_{Mc} =  \tfrac{(|B-C| - 1)^2}{B + C} \sim \chi^2_{1}.$$
\textbf{Example}:

\begin{center}
  %\renewcommand{\arraystretch}{1.5}
%\hspace*{\fill}
  \begin{tabular}{cc|cc}
      & & \multicolumn{2}{c}{Random Forest} \\
      & & correct & wrong \\
      \hline
      \multirow{2}{*}{Tree} & correct & 30 & 5 \\
      & wrong & 17 & 42 \\
  \end{tabular}
\end{center}

Calculate the test statistic:

$$\chi^2_{Mc} =  \frac{(|5-17| - 1)^2}{5 + 17} = 5.5 > 3.841 = \chi^2_{1,0.95}.$$

We can reject $H_0$ at a significance level of 0.05, i.e., we reject the hypothesis that the tree and the random forest have the same performance.
% and assume that our tree and our random forest classify the data differently

\end{frame}


\begin{frame}{Wilcoxon Signed Rank Test}
    \begin{itemize}
        \item The Wilcoxon Signed Rank Test ranks differences in performance of two classifiers, ignores the signs, and then compares ranks for positive and negative differences.
        \item In this case, $d_t$ denotes the difference between the performance scores of the two classifiers on the $t$-th data set:
        $$
        d_t = \widehat{\text{GE}}(\mathcal{I}_1, \mathcal{J}_t, \lamv_1, \rho) - \widehat{\text{GE}}(\mathcal{I}_2, \mathcal{J}_t, \lamv_2, \rho)
        $$
        \item Let $R^{+}$ denote the rank sum for data sets where algorithm B outperforms algorithm A and $R^{-}$ for the opposite:
        $$
        R^{+} = \sum_{d_i > 0} \text{rank}(d_i) + \frac{1}{2} \sum_{d_i = 0} \text{rank}(d_i) \quad\quad R^{-} = \sum_{d_i < 0} \text{rank}(d_i) + \frac{1}{2} \sum_{d_i = 0} \text{rank}(d_i) \quad\quad T = \min(R^{+}, R^{-})  
        $$
        \item The approximately normal $z$-score corresponds to:
        $$
        z = \frac{T - \frac{1}{4}N(N + 1)}{\sqrt{\frac{1}{24}N(N + 1)(2N + 1)}}
        $$
    \end{itemize}
\end{frame}

% \begin{frame}{McNemar Test}

% \begin{itemize}
% \item McNemar test is a non-parametric test that compares the errors of two classifiers.
% \item classifiers have to be trained and tested on the same dataset.
% % \item $H_0$: the two classifiers have the same error rates, \\ $H_1$: the error rates differ
% \item Goodness-of-fit measure: Compares observed counts with expected distribution under $H_0$.
% \item Relies on contingency matrix and is very similar to $\chi^2$-Test.
% \item McNemar test considers if observations are classified correctly or not and does not provide any quantification of these measures.
% \item According Dietterich(1998), the McNemar test should only be used if $B + C > 20$.
% \end{itemize}
% %
    
% \end{frame}

%  \begin{frame}{T-test vs. McNemar Test}

% Apart from t-test being parametric and the McNemar test being non-parametric, two tests have other significant differences. Because of its nature, it is possible to construct intervals for t-test however it is not possible to work with intervals for McNemar test which only works with nominal performance measures. McNemar test takes into account that whether observations are classified correctly or not and it does not provide any quantification of these measures.

% \lz

% According to the Dietterich(1998), one can proceed with McNemar test when the number of disagreement between two classifiers $c_{01}^{Mc} + c_{01}^{Mc} $ is greater than 20.

% \end{frame}


\begin{frame}{Comparing Multiple Models}

\begin{itemize}
    \setlength\itemsep{1em}
    \item So far, we have only been comparing the performance of two models.
    \item When comparing multiple models, we can make use of an omnibus test.
    \item Omnibus hypotheses to be tested are:
        \begin{itemize}
            \item $H_0:$ All algorithms are equivalent in their performance and hence their average ranks should be equal.
            \item $H_1:$ The average rank for at least one algorithm is different.
        \end{itemize}
    \item Omnibus tests are only useful to indicate that at least one model performs differently than the others. If that is the case, we need to do subsequent pairwise comparisons.
    \item When using pairwise comparisons, we need to do a correction for multiple hypothesis testing, e.g., a Bonferroni correction.
\end{itemize}
    
\end{frame}
% \begin{frame}{Cochran's Q Test? (not a slide)}

%     \begin{itemize}
%         \item Generalization of McNemar test to 3 or more models.
%         \item Not as good as Friedman or Nemeniy?
%     \end{itemize}
% \end{frame}

%\section{Compare Multiple Models Multiple Domains}


\begin{frame}{Friedman test}
% So far we have only compared 2 models / algorithms on one dataset. The \textbf{Friedman test} can be used to compare multiple classifiers on multiple datasets. The hypotheses to be tested are:
% \begin{itemize}
% \item $H_0:$ All algorithms are equivalent in their performance and hence their average ranks should be equal.
% \item $H_1:$ The average rank for at least one algorithm is different.
% \end{itemize}

Suppose we want to evaluate $m$ data sets and $k$ algorithms. The Friedman test first ranks algorithms for each data set separately:

\begin{itemize}
  \item Separately rank each algorithm for each dataset from best-performing (rank 1) to worst-performing, using any performance measure of interest.
  \item If there is a $d$-way tie after rank $r$, assign a rank of $ \left[(r+1) + (r+2) + ... + (r+d)\right] /d $ to all tied classifiers.
  \item $R_{ij}$ is the rank of algorithm $j$ on data set $i$.
\end{itemize}

After ranking algorithms, calculate the following quantities:

\begin{itemize}
  \item The overall mean rank:
  $ \bar{R} = \frac{1}{mk} \sum_{i=1}^{m} \sum_{j=1}^{k} R_{ij}. $
  \item The total sum of squares:
  $ SS_{Total} = m \sum_{j=1}^{k} (\bar{R}_{.j} - \bar{R})^2 $, where $\bar{R}_{.j} =  \frac{1}{m} \sum_{i=1}^{m} R_{ij}$.
  \item The error sum of squares:
  $ SS_{Error} = \frac{1}{m(k-1)} \sum_{i=1}^{m} \sum_{j=1}^{k} (R_{ij} - \bar{R})^2. $
\end{itemize}

The Friedman test statistic is calculated as:

$${\chi_F}^2 = \frac{SS_{Total}}{SS_{Error}} \sim \chi_{k-1}^2 \text{ for large \textit{m} ($>15$) and \textit{k} ($>5$).}$$
\end{frame}



% \begin{frame}{Friedman test}
% After obtaining the rank for each algorithm $j$ on different datasets $i$, calculate the following quantities:

% \begin{itemize}
%   \item The overall mean rank:
%   $ \bar{R} = \frac{1}{mk} \sum_{i=1}^{m} \sum_{j=1}^{k} R_{ij}. $
%   \item The total sum of squares:
%   $ SS_{Total} = m \sum_{j=1}^{k} (\bar{R}_{.j} - \bar{R})^2 $, where $\bar{R}_{.j} =  \frac{1}{m} \sum_{i=1}^{m} R_{ij}$.
%   \item The error sum of squares:
%   $ SS_{Error} = \frac{1}{m(k-1)} \sum_{i=1}^{m} \sum_{j=1}^{k} (R_{ij} - \bar{R})^2. $
% \end{itemize}

% % \lz

% The Friedman test statistic is calculated as:

% $${\chi_F}^2 = \frac{SS_{Total}}{SS_{Error}} \sim \chi_{k-1}^2 \text{ for large \textit{m} ($>15$) and \textit{k} ($>5$).}$$

% In the case of smaller $m$ and $k$, the $\chi^2$ approximation is imprecise and a look-up of $\chi_F^2$ values that were approximated specifically for the Friedman test is suggested.
% \end{frame}


\begin{frame}{Post-hoc tests}
A Friedman test checks if all algorithms are ranked equally or not. However, it does not provide information w.r.t. the best-performing algorithm.
To address this issue, post-hoc tests can be used.

\lz \textbf{Post-hoc Nemenyi test}:
\begin{itemize}
\item Compares all algorithms pairwise to find the best-performing algorithm after $H_0$ of the Friedman test was rejected.
\item For $m$ data sets and $k$ algorithms, $\frac{k(k-1)}{2}$ comparisons are made.
\item Calculates the average rank of algorithm $j$ on all $m$ data sets: $\bar{R}_{.j} =\frac{1}{m} \sum_{i=1}^m R_{ij}.$
\end{itemize}

\lz %The performance of two algorithms are significantly different when their average ranks differ by at least the significance level $q_{\alpha}$.
For any two algorithms $j_1$ and $j_2$, we compute the test statistic as:
$$q = \frac{\bar{R}_{.j_1} - \bar{R}_{.j_2}}{\sqrt{\frac{k(k+1)}{6m}}}.$$
\end{frame}

\begin{frame}{Post-hoc tests}
    
\textbf{Post-hoc Bonferroni-Dunn test}:
    
    \begin{itemize}
    \item Compares all algorithms with a baseline (i.e., $k-1$ comparisons).
    \item It is used after a Friedman test to find which algorithms differ from the baseline significantly.
    \item It uses the Bonferroni correction to prevent randomly accepting one of the algorithms as significant due to multiple testing.
    \end{itemize}
    The test statistic is the same as before:
    $$q = \frac{\bar{R}_{.j_1} - \bar{R}_{.j_{baseline}}}{\sqrt{\frac{k(k+1)}{6m}}}.$$
    
    The performances of $j_1$ and $j_2$ differ significantly if $|q| > q_{\alpha}$, where the critical value $q_{\alpha}$ is obtained from a table of the studentized range statistic, scaled through division by $\sqrt{2}$.
    %Another difference between Bonferroni-Dunn Test with Nemenyi test is that total number of comparisons are 1/k of Nemenyi test.
\end{frame}
\begin{frame}{Guidelines on Hypothesis Tests}

\begin{itemize}
    \item There is no golden standard for using hypothesis testing in performance evaluation. Often, tests are based on questionable or unverifiable assumptions or conclusions.
    \item The Wilcoxon and Friedman tests have been demonstrated to be most useful CITE. They partially fulfil commensurability, do not assume normal distributions or homogeneity of variance, and can be applied to any evaluation metric (e.g., accuracy, error ratios).
    \item For comparisons of two models, use the Wilcoxon test. For comparisons of multiple models, use the Friedman test.
    \item  "Tests provide certain reassurance about the validity and non-randomness of the published
results. For that to be true, they should be performed correctly and the resulting conclusions should
be drawn cautiously. On the other hand, statistical tests should not be the deciding factor for or
against publishing the work. Other merits of the proposed algorithm that are beyond the grasp of
statistical testing should also be considered and possibly even favoured over pure improvements in
predictive power." (Japkowiecz)
\end{itemize}
    
\end{frame}

xx
\endlecture
\end{document}
